% Autorska priprava k Realisticky.cz 020310–020311; ulohy jsou autorske a obmenene.
% Zamerne nenahrazuje vyklad presnymi kopiemi cizich textu/prikladu.
\documentclass[11pt,a4paper]{article}
\input{../../sablony/pripravy/priprava.sty}
\usepackage{tikz}
\nastavPripravu{1. ročník}{Rovnice a nerovnice}{81}{Soustavy: volba metody a procvičení}
\begin{document}
\blok{Návaznost a motivace}
Obě metody jsou známé. Žáci budou rozhodovat, kdy je výhodnější dosazovat a kdy sčítat. Cílem není nový pojem, ale jistota v samostatném postupu.
\blok{Princip a odvození}
Nejdřív rychle posuď koeficienty: obsahuje-li některá rovnice samostatnou neznámou s koeficientem $1$ nebo $-1$, bývá výhodné dosazování. Mají-li členy stejné či opačné koeficienty, bývá rychlejší sčítání. Vždy ověř obě rovnice.
\blok{Řešené příklady: od jednoduchého ke složitějšímu}
\begin{prikladReseniBox}[title={Příklad 1}]
Zvolte metodu a vyřešte $\begin{cases}4x+y=17,\\x-y=3.\end{cases}$
\tcblower
\[
\begin{aligned}
5x&=20 \\
 x&=4 \\
 y&=1
\end{aligned}
\]

Sčítání je velmi krátké. Oba členy s $y$ se okamžitě vyruší.
\[
\begin{aligned}
K &= \{[4;1]\}
\end{aligned}
\]
\end{prikladReseniBox}
\begin{prikladReseniBox}[title={Příklad 2}]
Zvolte metodu a vyřešte $\begin{cases}2x+3y=13,\\x+y=5.\end{cases}$
\tcblower
\[
\begin{aligned}
x&=5-y \\
 2(5-y)+3y&=13 \\
 y&=3,\quad x=2
\end{aligned}
\]

Zde se dobře dosazuje z druhé rovnice.
\[
\begin{aligned}
K &= \{[2;3]\}
\end{aligned}
\]
\end{prikladReseniBox}
\poznamka{Procvičení je vhodné zadat s volbou metody: ke každé vyřešené úloze má žák napsat jednou větou, proč metodu zvolil.}
\blok{Samostatné procvičování}
\begin{zadaniReseniBox}[title={Úloha 1}]
$\begin{cases}x+y=9,\\2x-y=3.\end{cases}$
\tcblower
\[
\begin{aligned}
3x&=12 \\
 x&=4,\quad y=5 \\[2mm]
K &= \{[4;5]\}
\end{aligned}
\]
\end{zadaniReseniBox}
\begin{zadaniReseniBox}[title={Úloha 2}]
$\begin{cases}4x-y=9,\\x+y=6.\end{cases}$
\tcblower
\[
\begin{aligned}
5x&=15 \\
 x&=3,\quad y=3 \\[2mm]
K &= \{[3;3]\}
\end{aligned}
\]
\end{zadaniReseniBox}
\begin{zadaniReseniBox}[title={Úloha 3}]
$\begin{cases}2x+2y=8,\\x+y=5.\end{cases}$
\tcblower
\[
\begin{aligned}
2(x+y)&=8\Rightarrow x+y=4 \\
 x+y&=5 \\[2mm]
K &= \varnothing
\end{aligned}
\]
\end{zadaniReseniBox}
\blok{Volitelná rezerva}
\begin{zadaniReseniBox}[title={Úloha 1}]
Jsou $\begin{cases}3x-2y=5,\\6x-4y=10\end{cases}$ dvě nezávislé podmínky?
\tcblower
\[
\begin{aligned}
6x-4y&=2(3x-2y) \\
 0&=0 \\
 y&=\frac{3x-5}{2}
\end{aligned}
\]

Ne, druhá rovnice nepřidává novou podmínku. Je nekonečně mnoho řešení.
\[
\begin{aligned}
K &= \left\{\left[t;\frac{3t-5}{2}\right]\mid t\in\mathbb R\right\}
\end{aligned}
\]
\end{zadaniReseniBox}
\blok{Domácí úkol}
\begin{zadaniReseniBox}[title={Domácí příprava}]
\begin{enumerate}[label=\alph*)]
\item
$\begin{cases}3x+y=10,\\x-y=2.\end{cases}$
\item
$\begin{cases}2x-y=4,\\x+2y=7.\end{cases}$
\item
$\begin{cases}4x+2y=12,\\2x+y=6.\end{cases}$
\item
$\begin{cases}5x-y=13,\\2x+y=8,\end{cases}$.
\item
$\begin{cases}2x+4y=6,\\x+2y=4,\end{cases}$.
\item
$\begin{cases}3x-6y=9,\\2x-4y=6,\end{cases}$.
\end{enumerate}
\tcblower
\begin{enumerate}[label=\alph*)]
\item
\[
\begin{aligned}
4x&=12 \\
 x&=3,\quad y=1 \\[2mm]
K &= \{[3;1]\}
\end{aligned}
\]
\item
\[
\begin{aligned}
y&=2x-4 \\
 5x&=15 \\
 x&=3,\quad y=2 \\[2mm]
K &= \{[3;2]\}
\end{aligned}
\]
\item
\[
\begin{aligned}
2(2x+y)&=12 \\
 0&=0 \\[2mm]
K &= \{[3-\frac t2;t]\mid t\in\mathbb R\}
\end{aligned}
\]
\item
\[
\begin{aligned}
7x&=21\\x&=3,\quad y=2 \\[2mm]
K &= \{[3;2]\}
\end{aligned}
\]
\item
\[
\begin{aligned}
2(x+2y)&=8\ne6 \\[2mm]
K &= \varnothing
\end{aligned}
\]
\item
\[
\begin{aligned}
x-2y&=3\\y&=t,\quad x=3+2t \\[2mm]
K &= \{[3+2t;t]\mid t\in\mathbb R\}
\end{aligned}
\]
\end{enumerate}
\end{zadaniReseniBox}
\poznamka{Orientační časové rozvržení: 5 + 12 + 12 + 10 + 6 = 45 min. Přednost má pochopení jádra; rezerva není součástí časového součtu. Domácí úlohy jsou určeny na samostatnou práci, řešení zůstávají učiteli.}
\end{document}
